\section{Methodology}

\subsection{Latent Space Model Formulation}
We assume a network $G=(V, E)$ with $n$ nodes. Let $z_i \in \mathbb{R}^d$ denote the unobserved latent position of node $i$. The probability of a tie between node $i$ and node $j$, denoted $A_{ij}=1$, is modeled as conditionally independent given their latent positions:

\begin{equation}
    P(A_{ij} = 1 | z_i, z_j, \alpha) = \frac{1}{1 + \exp(-(\alpha - d(z_i, z_j)))}
\end{equation}

where $\alpha$ is an intercept parameter controlling the baseline density of the network, and $d(z_i, z_j)$ is the distance between the nodes in the latent space.

\subsection{Geometries and Metrics}
The crucial component of our investigation is the choice of the distance metric $d(\cdot, \cdot)$, which is determined by the curvature of the latent manifold.

\paragraph{Euclidean Space (Zero Curvature)}
For the baseline model, we employ the standard Euclidean distance ($L_2$ norm):
\begin{equation}
    d_E(z_i, z_j) = ||z_i - z_j||_2 = \sqrt{\sum_{k=1}^d (z_{ik} - z_{jk})^2}
\end{equation}

\paragraph{Hyperbolic Space (Negative Curvature)}
To model hierarchical structures, we utilize the Poincaré disk model $\mathbb{D}^d = \{z \in \mathbb{R}^d : ||z|| < 1\}$. The distance between two points $u, v \in \mathbb{D}^d$ is given by:
\begin{equation}
    d_H(u, v) = \text{arccosh}\left( 1 + 2 \frac{||u - v||^2}{(1 - ||u||^2)(1 - ||v||^2)} \right)
\end{equation}
This metric grows exponentially as points approach the boundary of the disk ($||z|| \to 1$), providing sufficient space to embed tree-like hierarchies with low distortion.

\subsection{Estimation}
We estimate the latent positions $Z = \{z_1, \dots, z_n\}$ by maximizing the log-likelihood of the observed graph structure. The objective function is:

\begin{equation}
    \mathcal{L}(Z, \alpha) = \sum_{i < j} \left[ A_{ij} \log(p_{ij}) + (1 - A_{ij}) \log(1 - p_{ij}) \right]
\end{equation}

where $p_{ij} = P(A_{ij}=1 | z_i, z_j)$. We perform this maximization using the L-BFGS-B algorithm to handle the box constraints required for the Poincaré disk (i.e., ensuring $||z_i|| < 1$).

\subsection{Signed Network Extension}
To accommodate signed interactions (friends vs. foes), we extend the distance model to the Balanced Inner-Product model proposed by Tang and Zhu (2025). We introduce latent polar variables $v_i \in \{-1, 1\}$ to capture the sign of the relationship:

\begin{equation}
    \text{logit}(P(S_{ij} = 1 | z_i, z_j, v_i, v_j)) \propto v_i v_j - \beta d(z_i, z_j)
\end{equation}

where $v_i v_j > 0$ implies a positive tie (balance) and $v_i v_j < 0$ implies a negative tie (imbalance). The parameter $\beta$ controls the relative weight of geometric distance versus structural balance. This formulation allows us to model both the spatial proximity (via distance) and the signed relationship structure (via the inner product $v_i v_j$) simultaneously.

\subsection{Simulation Study Design}
Before applying the model to the Slashdot dataset, we validate our implementation via a simulation study.
\begin{enumerate}
    \item \textbf{Data Generation:} We simulate latent positions $Z_{true}$ from a known distribution (Gaussian for Euclidean, Uniform on Disk for Hyperbolic).
    \item \textbf{Network Formation:} Edges are generated using the probability model in Eq. (1).
    \item \textbf{Recovery:} We initialize the model with random noise and optimize Eq. (4) to recover $Z_{est}$.
    \item \textbf{Evaluation:} We assess performance by calculating the Area Under the ROC Curve (AUC) on a held-out set of edges.
\end{enumerate}

